\chapter{Les différents types de déchets}

L’objectif de ce premier module est de définir le plus clairement possible les termes décrivant les différentes catégories de déchets produits en France, de situer ces catégories les unes par rapport aux autres, de préciser les quantités en jeu et les moyens de collecte et de traitement aujourd’hui mis en œuvre. Ce module a également pour but de vous faire comprendre quelle logique sous-tend la réglementation des déchets en France, quels sont ses objectifs et ses moyens.


\section{Quels types de déchets ?}

Dans un premier temps, ce chapitre explique la plupart des termes utilisés pour distinguer les différents types de déchets produits par les communes, les industriels, le monde agricole, et les positionne les uns par rapport aux autres. Puis, dans un second temps, il donne des indications sur les quantités en jeu pour chaque type de déchets.

\subsection{Typologie}

Les différents types de déchets sont présentés dans le schéma ci-dessous.


\begin{figure}[H]
\centering
\includegraphics[width=0.7\textwidth]{Images/M1/panel_HW.PNG}
\end{figure}

\subsection{Mieux connaître les termes}

\textbf{Ordures ménagères (OM)}
\vspace{0.3cm}

Déchets issus de l’activité domestique quotidienne du foyer.
Dans certains cas, une collecte sélective est organisée pour séparer certains matériaux des OM :
\begin{itemize}
\item collecte des recyclables : bouteilles de verre, journaux, emballages en carton, boîtes de conserve et canettes, bouteilles en plastique;
\item collecte de fermentescibles : déchets de jardin (gazon, feuilles…) avec, parfois en plus, les déchets putrescibles de cuisine (épluchures, restes de repas…) ou même des papiers et cartons (ils pourrissent aussi…).
\end{itemize}
La commune est responsable de la collecte et de l’élimination des OM. Elles doivent être collectées au moins une fois par semaine.

\vspace{0.3cm}
\textbf{Flux de collecte sélective}
\vspace{0.3cm}

Un flux de collecte sélective est constitué par l’ensemble des matériaux collectés en même temps dans le même contenant. Selon le contexte, et en fonction notamment de l’unité de tri dont on dispose, on demandera aux habitants de séparer leurs matériaux recyclables en plusieurs flux :

Exemple

\begin{figure}[H]
\includegraphics[width=\textwidth]{Images/M1/selectivesorting.PNG}
\end{figure}


\vspace{0.3cm}
\textbf{Encombrants}
\vspace{0.3cm}

Déchets produits occasionnellement par un ménage qui, du fait de leur poids ou de leur volume (gravats, déchets de jardin, meubles usagés…), peuvent difficilement être collectés avec les OM.
\vspace{0.3cm}

La commune est responsable de leur collecte et de leur élimination. Elle peut proposer une collecte en porte-à-porte (une fois par mois ou une fois par trimestre) ou mettre à la disposition des particuliers une déchetterie.
Dans certains cas, les communes disposent des bennes sur les trottoirs pour récupérer les déchets encombrants.


\vspace {1cm}
\textbf {Déchets dangereux des ménages ou déchets ménagers spéciaux (DDM)}
\vspace {0.3cm}

Déchets dangereux pour la santé produits par les ménages et qui sont souvent liés aux activités de jardinage (traitements des végétaux, désherbants…), de bricolage (peintures, solvants, décapants, aérosols…) ou bien encore à l’entretien de la voiture (huiles de vidange, batteries, antigels, liquides de freinage…). On peut y ajouter les piles, les piles rechargeables, la plupart des produits d’entretien (déboucheurs d’évier, produits anti-tartre…), les insecticides, etc.
\vspace {0.3cm}

La commune est responsable de la collecte et de l’élimination de ces déchets. Ils ne doivent pas être collectés avec les OM. Elle peut proposer une collecte périodique, comme par exemple le service Planète (prestation Onyx) : une camionnette est stationnée sur les marchés une fois par mois et les particuliers peuvent y apporter leurs déchets dangereux. La collecte des déchets dangereux peut aussi être réalisée en déchetterie.

\vspace{0.3cm}
\textbf{Déchetterie}
\vspace{0.3cm}

Lieu clos et gardienné où les particuliers peuvent apporter leurs encombrants, leurs déchets de jardinage et de bricolage et, de plus en plus souvent, leurs déchets dangereux. Tous ces déchets sont déposés par le particulier dans des caissons ou dans des conteneurs spécifiques sous le contrôle du gardien. L’objectif de ce tri est de pouvoir recycler ou valoriser le maximum de matériaux.

\vspace{0.3cm} 
\textbf{Déchets de nettoiement}
\vspace{0.3cm}

Tous les déchets produits sur les espaces publics de la commune : déchets de balayage des rues, vidage des corbeilles, déchets des marchés.

\vspace{0.3cm}
\textbf{Déchets d’espaces verts}
\vspace{0.3cm}

Les déchets végétaux des parcs et jardins publics de la collectivité, les feuilles et les tailles des arbres plantés le long des voies publiques, les déchets végétaux générés par les parterres et les bacs de fleurs de la collectivité.

\vspace{0.3cm}
\textbf{Déchets de l’assainissement}
\vspace{0.3cm}

Boues résiduelles après dépollution des eaux usées (boues de station d’épuration des eaux usées), résidus de curage des égouts.

\vspace{0.3cm}
\textbf{Déchets industriels banals ou déchets industriels et commerciaux (DIB / DIC)}
\vspace{0.3cm}



Déchets issus de l’activité économique (entreprises ou établissements), dont la nature est proche de celle des déchets des ménages (papiers, cartons, plastiques, métaux, bois, déchets putrescibles…) et qui peuvent être traités dans les mêmes installations de traitement que les ordures ménagères.

\vspace{0.3cm}

Lorsque les quantités restent faibles, la collecte de ces déchets peut être réalisée en même temps que les ordures ménagères. C’est presque toujours le cas pour les petits commerces et artisans de centre ville, et pour les déchets des écoles, des administrations, des casernes. Cependant, la commune n’a aucune obligation de prendre en charge ces déchets qui n’ont pas été produits par des ménages, mais par des professionnels.

\vspace{0.3cm}

Quand les quantités sont importantes (usines, supermarchés…), les DIB/DIC sont collectés par des entreprises privées qui proposent au producteur des moyens de collecte adaptés aux gros volumes (caissons, compacteurs…). Cette collecte cherchera à séparer les différents déchets afin de favoriser la valorisation.
\vspace{0.3cm}

Remarque importante
Tout professionnel qui produit plus de 1100 litres d’emballages usagés par semaine doit s’organiser pour permettre la valorisation de ces emballages (recyclage ou valorisation énergétique).

\vspace{0.3cm}

La collecte et l’élimination de ces déchets sont de la responsabilité de l’entreprise ou de l’établissement producteur. Quand ces déchets sont collectés en même temps que les OM, cette responsabilité est transférée à la collectivité qui, en cas de problème, peut se retourner contre l’établissement s’il est identifié.
\vspace{0.3cm}

\vspace{0.3cm}
\textbf{Déchets dangereux (DD)}
\vspace{0.3cm}

Ce terme remplace l’ancienne désignation « Déchets Industriels Spéciaux ». Il s’agit de déchets produits par des entreprises et qui, de par leur caractère polluant, doivent être éliminés dans des installations spécialisées.
Ces déchets peuvent être :

\begin{itemize}
\item liquides : organiques (hydrocarbures, solvants, etc.), minéraux (acides, bases, etc.) ;
\item boues : boues de cabine de peinture, boues d'hydroxydes métalliques provenant de traitements de surface ;
\item solides : goudrons, terres contaminées, catalyseurs, produits polymérisés, résidus solides de distillation, résidus d'épuration des fumées d'incinérateurs, amiante, etc.
\end{itemize}

Un suivi très strict de ces déchets est réalisé par le ministère de l’Environnement.
La collecte et l’élimination de ces déchets sont de la responsabilité de l’entreprise ou de l’établissement producteur.

\vspace{0.3cm}

\textbf{Déchets dangereux diffus (DDD)}
\vspace{0.3cm}

\begin{figure}[H]
\centering
\includegraphics[width=0.5\textwidth]{Images/M1/DHW.PNG}
\end{figure} 

Ils sont de même nature que les déchets dangereux (DD), mais sont produits en faibles quantités par des petites entreprises ou par les laboratoires des lycées ou des universités. De ce fait, les moyens de collecte sont différents : par exemple, SARPI-Onyx propose le service Coralie, basé sur la mise à disposition de bonbonnes spécifiques à chaque produit chimique, collectées périodiquement.
La collecte et l’élimination de ces déchets sont de la responsabilité de l’entreprise ou de l’établissement producteur.


\vspace {0.3cm}
\textbf{Déchets inertes}
\vspace {0.3cm}

Déchets non susceptibles d’évolution physico-chimique : qui ne brûlent pas, qui ne fermentent pas, qui ne diffusent pas d’éléments polluants dans l’eau.
Deux grands types d’inertes :
\begin{itemize}
\item résidus extraits des mines (75\%) ;
 \item certains déchets de démolition (25\%), sauf bois, charpentes, faux plafonds, isolants, plaques de plâtre, etc.
\end{itemize}

Remarque : les déchets végétaux ne peuvent en aucun cas être considérés comme des inertes et être éliminés dans les mêmes conditions que ceux-ci.

La collecte et l’élimination des déchets inertes sont de la responsabilité de l’entreprise ou de l’établissement producteur.

\vspace {0.3cm}
\textbf{Déchets d’activités de soins à risque infectieux (DASRI)} 
\vspace{0.3cm}

Ils regroupent les déchets produits par les hôpitaux et les cliniques (déchets hospitaliers), mais aussi par tous les autres professionnels de la santé : médecins généralistes et spécialistes, dentistes, infirmières libérales, vétérinaires, laboratoires d’analyses médicales, etc.
On les classe en deux catégories :


\begin{itemize}
\item Les déchets contaminés, susceptibles de transmettre des germes ou maladies : déchets anatomiques, pansements, seringues… (tous les déchets pour les services de maladies infectieuses). Ils peuvent être produits à l’hôpital, mais aussi dans le cabinet du médecin, du dentiste ou même à domicile (diabétiques, par exemple). Dans ces deux derniers cas, on parlera de déchets contaminés diffus.

\item Les déchets assimilables aux ordures ménagères : reliefs de repas, cartons, déchets de bureaux… qui, comme les DIB/DIC, sont collectés avec les ordures ménagères ou par une entreprise privée.
Comme pour les DIB/DIC, si l’établissement produit plus de 1100 litres d’emballages par semaine, il doit organiser la séparation de ces emballages des autres déchets en vue de leur valorisation. La collecte et l’élimination des déchets d’activité de soins sont de la responsabilité de l’établissement hospitalier ou du professionnel de la santé qui les a produits.
Quand ces déchets sont collectés en même temps que les OM (fraction assimilable aux OM), cette responsabilité est transférée à la collectivité qui, en cas de problème, peut se retourner contre l’établissement ou le producteur s’il est identifié.
\end{itemize}



\subsection{Les déchets en chiffres}

\begin{minipage}{0.65\textwidth}
    \vspace{0.3 cm}
    Plus de 600 millions de tonnes de déchets par an, dont :
    \vspace{0.3 cm}
    
    \textbf{Déchets agricoles}
    420 millions de tonnes/an
    
    \begin{enumerate}
    \item Élevage (fumier solide ou liquide, etc.) : 280 millions de tonnes/an
    \item Production agricole : 55 millions de tonnes/an
    \item Foresterie : 40 millions de tonnes/ an
    \item Industrie alimentaire : 45 millions de tonnes/an
    \end{enumerate}
\end{minipage}%
\begin{minipage}{0.35\textwidth}
    \begin{figure}[H]
    \includegraphics[width=\textwidth]{Images/M1/farmingwaste.PNG}
    \end{figure}
\end{minipage}


\vspace{0.3 cm}


\vspace{1 cm}
\textbf{Déchets industriels}
\vspace{0.3cm}

\begin{minipage}{0.35\textwidth}
    \begin{figure}[H]
        \includegraphics[width=\textwidth]{Images/M1/industrialwaste.PNG}
    \end{figure}
\end{minipage}%
\begin{minipage}{0.65\textwidth}
    150 millions de tonnes/an

    \begin{enumerate}
        \item Déchets inertes : 100 millions de tonnes/an
        \item Déchets industriels ordinaires (DIB) : 40 millions de tonnes/an
        \item Déchets dangereux (HW) : 7 millions de tonnes/an	
    \end{enumerate}
\end{minipage}









\vspace{1 cm}
\textbf{Déchets d’activités de soins}
\vspace{0,3 cm}

0,8 million de tonnes/an

2 à 6 kg par lit et par jour
	Déchets hospitaliers contaminés : 0,21 million de tonnes/an
	Déchets hospitaliers considérés comme déchets ménagers : 0,49 million de tonnes/an
	Déchets de santé contaminés diffus : 0,1 million de tonnes/an

\vspace{0.3cm}
\textbf{Déchets municipaux}
\vspace{0.3cm}

\begin{minipage}{0.7\textwidth}
    53 millions de tonnes/an :
    \begin{enumerate}
    \item Déchets ménagers : 19,2 millions de tonnes/an
    \item Collectes sélectives : 1,6 million de tonnes/an
    \item Déchets des commerçants/entreprises collectés avec ho déchets ménagers : 5,2 millions de tonnes/an soit 434 kg par habitant et par an
    \item Déchets ménagers occasionnels 4,5 millions de tonnes/an soit 80 kg par habitant et par an
    \item Déchets d'assainissement et d'assainissement : 22,5 millions de tonnes/an soit 380 kg par habitant et par an
    \end{enumerate}
\end{minipage}
\begin{minipage}{0.3\textwidth}
\begin{figure}[H]
\includegraphics[width=\textwidth]{Images/M1/municipalwaste.PNG}
\end{figure}
\end{minipage}

%\begin{figure}[H]
%\includegraphics[width=\textwidth]{Images/M1/currentwastecompo.PNG}
%\end{figure}

\begin{figure}[H]
    \centering
    \includegraphics[width=\textwidth]{Images/M1/currentwastecompo.PNG}
\end{figure}





\section{Quels traitements pour quels déchets ?}

Ce chapitre présente les moyens de traitement mis en œuvre pour les différents types de déchets préalablement définis et les perspectives d’évolution de ces techniques.

\subsection{Les déchets agricoles}

La plupart des déchets agricoles sont réutilisés ou valorisés :
\begin{itemize}
\item par épandage : fumier solide ou liquide, eau de déshydratation de luzerne, etc. ;
\item par valorisation énergétique : paille, écorce, sciure de bois ;
\item comme aliment pour animaux ou pour humains ;
\item dans l'industrie cosmétique;
\item dans l'industrie pharmaceutique.
\end{itemize}

\vspace{0.3cm}
\textbf{Perspectives d’évolution :}
\vspace{0.3cm}

\begin{itemize}
\item des efforts à poursuivre sur les productions polluantes (tannerie, mégisserie : travail du cuir comme teinture…, graisses…);
\item le problème des farines de viandes est en cours de résolution : leur incinération étant maintenant obligatoire;
\item les débouchés en alimentation animale sont fragiles car ces produits sont en concurrence avec les céréales dont les cours sont en baisse continue.
\end{itemize}

\subsection{Les déchets industriels}

\textbf{Déchets inertes}
\vspace{0.3cm}

Ces déchets sont actuellement :
\begin{itemize}
\item utilisés comme décharge (Sites de stockage définitif - CSDU) ;
\item récupéré en remblais de carrières ;
\item recyclés sous forme de granulats (déchets inertes broyés réutilisés pour fabriquer du béton, etc.).
\end{itemize}

\textbf{Perspectives de changement :}
\vspace{0.3cm}

\begin{itemize}
\item séparer autant que possible les vrais déchets inertes des autres déchets de démolition (voir définition des « déchets inertes ») ;
\item les déchets de démolition sont de plus en plus recyclés sous forme de granulés.
\end{itemize}

\vspace{0.3cm}
\textbf{Déchets industriels banals} 
\vspace{0.3cm}

Ces déchets, de nature très proche des ordures ménagères, suivent en général les mêmes voies de traitement que celles-ci, ou des voies relativement proches. Cependant :

\begin{itemize}

\item le recyclage des DIB est bien plus développé que pour les OM, et ceci depuis longtemps : cartons, palettes, métaux, housses plastiques… Le recyclage est facilité par la présence de quantités importantes en un même point, ainsi que par des gisements propres et de bonne qualité;

\item la part de la mise en décharge est encore importante;
\item en 1993, on ne connaissait pas l’exutoire précis de près d’un tiers des DIB (brûlage au fond de la cour de nombreuses entreprises, par exemple).
\end{itemize}
\vspace{0.5cm}


\vspace{1cm}
\textbf{Perspectives de changement :}
\vspace{0.3cm}

\begin{itemize}

\item réduire les quantités de déchets en modifiant les lignes de production ;
\item les déchets industriels doivent être valorisés si l'entreprise produit plus de 1100 litres par semaine.
\end{itemize}

\vspace{0.3cm}
\textbf{Déchets dangereux}
\vspace{0.3cm}

De par leur nature, ces déchets doivent suivre des filières spécialisées de traitement :
\begin{itemize}

\item déchets organiques (huiles, solvants, boues de peinture, etc.) : incinérés dans un centre spécialisé ;
\item déchets minéraux liquides (acides, bases, bains de chrome, etc.) : traitement physico-chimique pour neutraliser ou concentrer et extraire les polluants ;
\item déchets solides (goudrons, cires, terres contaminées, produits polymérisés, résidus solides de distillation, catalyseurs, sels de trempe au cyanure, résidus d'épuration des fumées d'incinération, etc.) : incinération dans des centres spécialisés, et inertage des déchets ultimes pour fixer les polluants (les rendre insoluble dans l'eau), suivi d'un stockage en site de stockage de déchets résiduels (CSDU).
\end{itemize}

\vspace {0.3cm}
\textbf{Quelques chiffres (1993) :}
\vspace{0.3cm}

Valorisation = 32 du tonnage d'Eau
Traitement = 17\% du tonnage d'Eau
Stockage = 51\% du tonnage d'Eau
50\% du tonnage d'Eau est traité dans les unités de traitement internes par le entreprises génératrices de déchets.
\vspace{0.3cm}

\textbf{Quelques prix pour l'élimination du HW ?}
\vspace{0.3cm}

Traitement par incinération : 100 - 105 € / tonne
Traitement physico-chimique : 80 - 150 € / tonne
Stabilisation des déchets ultimes et stockage en CSDU : 300 € / tonne (hors taxes appliquées par l'Agence de l'Environnement et de la Maîtrise de l'Energie - Ademe)

\vspace {0.3cm}
\textbf{Perspectives de changement :}
\vspace{0.3cm}

\begin{itemize}
\item réduire les quantités de déchets en aménageant l’outil de production ;
\item plus grande sélectivité des déchets admis en centre de stockage qui est réservé, depuis le 1er juillet 2002, aux résidus ultimes “ stabilisés ”. Le CSDU remplace les décharges de classe 1 et 2.
\end{itemize}

\subsection{Les déchets d’activités de soins contaminé}

\vspace{0.3cm}
\textbf{Précautions sanitaires clés} :
\vspace{0.3cm}

\begin{itemize}
\item garantir un système de collecte entièrement hermétique pour éviter la propagation des germes pathogènes contenus dans ces déchets ;
\item détruit les germes pathogènes par incinération, chauffage ou désinfection chimique.
\end{itemize}

Solutions de traitement multiples :
\begin{itemize}

\item incinération en centre spécialisé de traitement des DASRI contaminés;
\item incinération dans des usines d’incinération d’ordures ménagères, si un système de vidage automatique des déchets dans la trémie du four est en place : aucun opérateur ne doit avoir de contact avec les DASRI, pas de vidage dans la fosse destinée aux OM. Ces DASDI ne doivent pas représenter plus de 10\% des tonnages incinérés par l’usine;
\item décontamination par chauffage (micro-ondes, vapeur…) ou désinfection chimique, puis traitement des déchets désinfectés ou “ banalisés ” comme des ordures ménagères classiques par incinération.
\end{itemize}

\vspace{0.3cm} 
\textbf{Déchets d’activités de soins non contaminés}
\vspace{0.3cm}

Ceux-ci sont considérés comme des déchets industriels banals (déchets de professionnels), ils sont bien souvent collectés en même temps que les ordures ménagères et suivent les mêmes voies de traitement.
On se rend bien compte ici qu’il est indispensable que l’hôpital maîtrise parfaitement les circuits de collecte de ses déchets contaminés et non contaminés…


Perspectives de changement :
\begin{itemize}
\item disparition des incinérateurs internes aux hôpitaux (trop petits pour être bien gérés);
\item comme toute entreprise, les hôpitaux sont tenus de faire valoriser les emballages qu’ils produisent s’ils en produisent plus de 1100 litres par semaine;
\item développement indispensable de la collecte des déchets d’activités de soins, produits de manière diffuse par les professions libérales de la santé.
\end{itemize}

\subsection{Les ordures ménagères en France}

\textbf{Le traitement des OM aujourd’hui}
\vspace{0.3cm}

En 1993, près de 60\% des déchets ménagers étaient stockés en décharge ou traités sans aucune forme de valorisation.

\begin{figure}[H]
\includegraphics[width=\textwidth]{Images/M1/typehouseholdtreatment.PNG}
\end{figure}

\textbf{Filière globale de traitement}
\vspace{0.3cm}

La combinaison de plusieurs moyens de traitement (recyclage, compostage, incinération) permet de réduire la quantité de résidus ultimes produite.


\begin{figure}[H]
\includegraphics[width=\textwidth]{Images/M1/oneton.PNG}
\end{figure}

Cette combinaison du recyclage, du compostage et de l’incinération ne peut pas être 
généralisée sur tout le territoire national. Dans les départements très ruraux, les tonnages sont insuffisants pour qu’une unité d’incinération soit mise en place avec des coûts de traitement raisonnables. Dans ces départements, l’organisation du traitement des OM pourra s’appuyer sur le recyclage des emballages et des journaux-magazines, et sur le compostage de la fraction organique, les autres résidus étant stockés en décharge.

\vspace{0.3cm} 

La notion de résidus ultimes s’adapte au contexte local et tient compte des coûts.

\section{La réglementation}

Il s’agit ici de comprendre la logique et les objectifs de la réglementation sur les déchets en France.


\textbf{Le bilan peu satisfaisant à la fin des années 1980}
\vspace{0.3cm}

\begin{itemize}
\item 60\% des déchets municipaux et industriels étaient alors enfouis en décharge ; 
\item  le recyclage et la valorisation restaient marginaux ;
\item l’ouverture de nouveaux sites de stockage des déchets était de plus en plus difficile.
\end{itemize}

Gestion des déchets basée essentiellement sur le stockage en décharge + rejet par la population de ce système de traitement = à court terme, un problème!

\vspace{0.3cm}

\textbf{Mise en place d’une nouvelle politique de gestion des déchets}
\vspace{0.3cm}

Un thème majeur : limiter l'impact des centres de stockage
Pour cela :

\begin{itemize}
\item Produire moins de déchets : au niveau de la fabrication des produits ou lors de leur utilisation.
			Label NF Environnement (peintures, piles…)

\item Réduire les quantités à stocker en décharge en favorisant le recyclage et la valorisation énergétique (chaleur ou électricité) ou organique (compost).

    \begin{itemize}
    \item Donner la priorité au recyclage et à la valorisation : objectif de valorisation par recyclage ou par compostage de 50\% des déchets ménagers. L’application directe commence par les emballages (voir ci-dessous).
    \item  Réserver les centres de stockage aux déchets ultimes à partir du 1er juillet 2002.
    \end{itemize}
\end{itemize}

\vspace{0.3cm}
\textbf{Définition}
\vspace{0.3cm}

Déchet ultime : au sens strict, c’est un résidu du traitement des déchets dont on ne peut plus extraire ni élément recyclable, ni matière organique, ni énergie.
Ce sont, par exemple, des résidus de dépollution des fumées d’incinération. Cette notion est à moduler en fonction du temps : un déchet ultime aujourd’hui n’en sera peut-être plus un demain, car la technique permettra de valoriser encore le produit. D’autre part, cette notion est à moduler en fonction du contexte local : les refus de compostage qui, techniquement, peuvent être valorisés énergétiquement, pourront être considérés comme déchets ultimes pour une collectivité rurale, car leur incinération engendrerait des coûts très élevés, ce qui n’est pas le cas en zone urbaine où la population et les tonnages sont plus importants, donc les coûts moins élevés.

\vspace{0.3cm}
\textbf{Deux outils administratifs pour limiter les quantités stockées en décharge :}
\vspace{0.3cm}

\begin{itemize}

\item Taxe sur le stockage : en 1998 une taxe de 40 F/t est prélevée sur toutes les tonnes de déchets ménagers stockées en décharge (en 1999, cette taxe est prévue d’atteindre à 60 F/t). Ces fonds permettent à l’Ademe (voir chapitre “ Les partenaires ”) de subventionner les collectivités lors de la mise en place de collectes sélectives ou d’unités de tri et de traitement.
\item Les plans départementaux : tous les départements doivent avoir rédigé un plan départemental de gestion des déchets ménagers et assimilés. Ce plan doit prévoir l’organisation de la collecte et du traitement, fixer des objectifs de valorisation, fixer des zones de collecte pour chaque équipement de traitement existant ou à créer.
\end{itemize}

Améliorer les techniques de stockage :
\begin{itemize}

 \item Règles d’admissibilité des déchets en centre de stockage : séparer les résidus ultimes, des déchets pouvant être encore valorisés (reprise éventuelle de ceux-ci plus tard?).

\item Chaque alvéole doit être complètement isolée du milieu extérieur : tous les lixiviats (rejets liquides recueillis au fond des alvéoles) et biogaz doivent être récupérés et traités.

\item Mise en place d’un suivi à long terme : pendant trente ans après la fin de l’exploitation du centre de stockage, l’exploitant reste responsable de sa surveillance et de son suivi. 
Durant toute cette période, les jus et biogaz doivent être captés et traités, et la bonne étanchéité des cellules surveillée.

\end{itemize}
 
\section{Quelques procédés industriels générateurs de DID}


En 2002, l’ADEME estimait la production française de déchets dangereux à environ 6 millions de tonnes. Les secteurs d’activités à l’origine de ces déchets semblent très divers (cf. figure n°1). Néanmoins, deux familles de secteurs d’activité cohérents contribuent de façon majoritaire au bilan total :
\begin{itemize}
\item Le secteur Métallurgique : Déchets mécaniques, électriques et électroniques : 35\% des quantités totales.
\item Le Secteur Chimique : Pétrochimie, Chimie Fondamentale, Chimie Fine, Pharmaceutique, Parachimie : 31\% des quantités totales.

\end{itemize}

\begin{figure}[H]
\centering
\includegraphics[width=0.5\textwidth]{Images/M1/fewindus.PNG}
\end{figure}